\section{Introduction}
\label{sec:introduction}

An agent is given one approval to perform a consequential action.
To improve reliability, its runtime forks two candidate plans and promises that only one will be retained.
Each branch prepares a different action under the approval.
This can be safe: the two promises are conditional on mutually exclusive futures, so one unit of authority can cover either outcome.
Now consider two lifecycle changes.
First, one branch sends its prepared request before selection.
The remote observer retains that effect even if the branch later loses.
Second, an operator merges both candidate results rather than choosing one.
In both cases, claims that were previously alternatives can now contribute to one durable history.
No branch-local capability check, workspace snapshot, or successful transaction by itself explains whether the change is authorized.

This boundary is becoming common.
Tool-using agents restore conversation checkpoints after errors, run best-of-\(N\) searches, delegate work to parallel subagents, resume old sessions after human edits, and merge artifacts from several branches.
The execution state being copied is heterogeneous: files, database rows, model context, plans, branch eligibility, approvals, capability ownership, revocation epochs, and remote effects can reside in different rollback domains.
Restoring all of it is often impossible; restoring only the workspace is insufficient.
The security-relevant question is not which bytes were copied, but which descendant outcomes may still become durable together and which effects have already escaped.

Existing foundations solve important parts of this problem.
Linear capabilities and consumable credentials prevent syntactic duplication from creating additional use-limited authority~\cite{bowers2007consumable,terauchi2006capability}.
State-continuity systems preserve monotone history across rollback~\cite{parno2011memoir,matetic2017rote,brandenburger2017lcm}.
Transactions and durable execution separate tentative work from settlement.
Recent agent systems provide branch isolation, first-commit-wins selection, explicit capabilities, semantic transactions, and commit-time checks~\cite{wang2026fork,zhang2026agentlibos,mohammadi2026atomix,chen2026cordon,santos2026committime}.
These mechanisms do not yet supply a common security contract for \emph{conditional authority promises} across a lifecycle that may later restore, expose, or merge their branches.

The tempting formulation is “forking must not fork authority.”
It is too coarse.
Always splitting authority at fork is safe but rejects useful late-bound exploration.
Keeping authority in a parent escrow and transferring it only after selection is also safe, and refutes any claim that topology is necessary for pure speculation.
Topology becomes security-relevant only when the runtime makes an advance commitment to a branch or lets an effect escape before the choice is durably resolved.
We therefore model a tentative claim as an enforceable promise: if its branch is retained and its grant remains live, the claim may pass the trusted commit gate without obtaining new authority.

Our central observation is that lifecycle operations change the \emph{support} of such promises.
A claim held by one alternative is conditional on futures containing that branch.
When its effect escapes, the claim becomes durable consumption that all remaining futures must carry.
A merge can similarly turn previously exclusive commitments into jointly retained commitments.
Restore-replace and restore-live may reconstruct identical local state, yet the former kills the old continuation while the latter makes it concurrent with the restored one.
Authorization must follow these lifecycle facts rather than checkpoint bytes.

We formalize this idea as \emph{authority continuity}.
For every set of branches that the trusted lifecycle permits to remain durable together, durable consumption plus their conditional claims must fit a typed grant.
The definition is deliberately simple.
Its content comes from connecting it to a concrete runtime semantics, proving why restored bytes omit necessary authorization facts, and deriving the weakest safe future restriction when conditional effects become irreversible.

\paragraph{Contributions.}
This paper makes four contributions.

\begin{enumerate}[leftmargin=*,itemsep=2pt]
  \item We give a concrete semantics for history-transforming execution that separates reconstructable agent state from a monotone security ledger and covers exclusive and parallel fork, checkpoint, replacing and live restore, selection, abort, escape, uncertain dispatch, merge, and revocation.
  We define an abstraction from this state into typed, uniquely identified conditional and durable claims over co-durable frontiers.

  \item We define authority continuity and prove abstraction soundness for conservative lifecycle contracts, with a converse only under exact frontiers and conjunctively realizable branch bundles.
  We prove that no snapshot-local checker can be both sound and maximally permissive, because replacing and live restore can present identical restored bytes but require opposite decisions.

  \item We identify \emph{escape promotion}, a lifecycle boundary missed by a topology-only account.
  A two-branch counterexample shows that moving a conditional claim into durable history can invalidate a safe exclusive choice.
  We derive its exact admission condition and the unique inclusion-largest safe restriction of the existing future family without new capacity or other cancellations.
  We prove atomic durable selection is a conservative instance and that batched promotion repairs commute.

  \item We give contextual rules for replacing and live restore, a policy-parametric merge rule, and an end-to-end preservation theorem over terminal claim IDs, durable consumption, and closed epochs.
  We connect the semantics to implementable representations without claiming classical mathematics as new.
  Structured choice/parallel contracts induce cographs and admit the standard linear cotree dynamic program; arbitrary pairwise conflict inherits independent-set hardness.
  We specify a small mechanization and executable mutation study, plus deterministic runtime litmus tests centered on the formal distinctions rather than model nondeterminism.
\end{enumerate}

The theory applies to adaptive non-LLM runtimes as well.
Agents make it urgent because they combine cheap semantic branching, long-lived delegated authority, dynamically synthesized effects, external tools, human resume, and artifact merge in one execution model.
